\documentclass[fleqn]{beamer}

\usepackage{booktabs}
\usepackage{subcaption}

\usepackage{stmaryrd}
\usepackage{mathpartir}
\usepackage{amssymb}
\usepackage{mathrsfs}
\usepackage{cmll}
\usepackage{xcolor}
\usepackage{makecell}
\usepackage{tikz-cd}
\usepackage{nccmath}
\usepackage{mathtools}
\usepackage{graphicx}

\usetikzlibrary{positioning}
% \usetikzlibrary{decorations.markings}
\usetikzlibrary{arrows}

\def\newelims{1}
\def\multnotation{1}
\ifx\newelims\undefined
  \def\newelims{0}
\fi
\ifx\multnotation\undefined
  \def\multnotation{0}
\fi

\newcommand{\fixme}[1]{{\color{red}\underline{\em{#1}}}}
\newcommand{\fixmeBA}[1]{{\color{red}{[Bob:{\em #1}]}}}


\newcommand{\bind}[2]{%
  \if\debruijn0%
  {#2}\{{#1}\}%
  \else%
  {#2}%
  \fi%
}
\newcommand{\ctx}[2]{%
  \if\multnotation0%
  #1^{\resctx{#2}}%
  \else%
  \resctx{#2}#1%
  \fi
}
\newcommand{\ctxvar}[3]{%
  \if\debruijn0%
  \if\multnotation0%
  #1 \stackrel{\rescomment{#3}}: #2%
  \else%
  \rescomment{#3}#1 : #2%
  \fi%
  \else%
  \rescomment{#3}#2%
  \fi%
}
\definecolor{res}{HTML}{008000}
\definecolor{resperm}{HTML}{008000}
\newcommand{\rescomment}[1]{{\color{res}#1}}
\newcommand{\rescommentperm}[1]{{\color{resperm}#1}}
\newcommand{\resctx}[1]{\rescomment{\mathcal{#1}}}
\newcommand{\resctxperm}[1]{\rescommentperm{\mathcal{#1}}}
\newcommand{\resmat}[1]{\rescomment{#1}}


\newcommand{\ann}[2]{#1 : #2}
\newcommand{\emb}[1]{\underline{#1}}

\newcommand{\base}[0]{\iota}

\newcommand{\fun}[2]{#1 \multimap #2}
\newcommand{\lam}[2]{%
  \if\debruijn0%
  \lambda #1.~#2%
  \else%
  \lambda #2%
  \fi%
}
\newcommand{\app}[2]{#1~#2}

\newcommand{\excl}[2]{\oc_{\rescomment{#1}} #2}
\newcommand{\bang}[1]{\left[#1\right]}
\newcommand{\bm}[4]{%
  \if\newelims0%
  \operatorname{bm}_{#1}(#2, \bind{#3}{#4})%
  \else%
  \if\debruijn0%
  \mathrm{let}~\bang{#3} = #2~\mathrm{in}~#4%
  \else%
  \mathrm{let}~\bang{-} = #2~\mathrm{in}~#4%
  \fi%
  \fi%
}

\newcommand{\tensorOne}[0]{1}
\newcommand{\unit}[0]{(\mathbin{_\otimes})}
\newcommand{\del}[3]{\if\newelims0%
\operatorname{del}_{#1}(#2, #3)%
\else%
\mathrm{let}~\unit = #2~\mathrm{in}~#3%
\fi}

\newcommand{\tensor}[2]{#1 \otimes #2}
\newcommand{\ten}[2]{(#1 \mathbin{_{\otimes}} #2)}
\newcommand{\prm}[5]{%
  \if\newelims0%
  \operatorname{pm}_{#1}(#2, \bind{#3, #4}{#5})%
  \else%
  \if\debruijn0%
  \mathrm{let}~\ten{#3}{#4} = #2~\mathrm{in}~#5%
  \else%
  \mathrm{let}~\ten{-}{-} = #2~\mathrm{in}~#5%
  \fi%
  \fi%
}

\newcommand{\withTOne}[0]{\top}
\newcommand{\eat}[0]{(\mathbin{_{\with}})}

\newcommand{\withT}[2]{#1 \with #2}
\newcommand{\wth}[2]{(#1 \mathbin{_\with} #2)}
\newcommand{\proj}[2]{\operatorname{proj}_{#1} #2}

\newcommand{\sumTZero}[0]{0}
\newcommand{\exf}[2]{\operatorname{ex-falso} #2}

\newcommand{\sumT}[2]{#1 \oplus #2}
\newcommand{\inj}[2]{\operatorname{inj}_{#1} #2}
\newcommand{\cse}[6]{%
  \if\newelims0%
  \operatorname{case}_{#1}(#2, \bind{#3}{#4}, \bind{#5}{#6})%
  \else%
  \if\debruijn0%
  \mathrm{case}~#2~\mathrm{of}~\inj{L}{#3} \mapsto #4
                            ~;~ \inj{R}{#5} \mapsto #6%
  \else%
  \mathrm{case}~#2~\mathrm{of}~\inj{L}{-} \mapsto #4 ~;~ \inj{R}{-} \mapsto #6
  \fi%
  \fi%
}

\newcommand{\listT}[1]{\operatorname{List} #1}
\newcommand{\nil}[0]{[]}
\newcommand{\cons}[2]{#1 :: #2}
\newcommand{\fold}[5]{\operatorname{fold} #1~#2~(#3,#4.~#5)}


\newcommand{\typed}[1]{\mathit{#1t}}
\newcommand{\resourced}[1]{\mathit{#1r}}


\newcommand{\lvar}{\mathrel{\mathrlap{\sqsupset}{\mathord{-}}}}
\newcommand{\sem}[1]{\left\llbracket #1 \right\rrbracket}

%\def\tobar{\mathrel{\mkern3mu  \vcenter{\hbox{$\scriptscriptstyle+$}}%
%                    \mkern-12mu{\to}}}
\newcommand\tobar{\mathrel{\ooalign{\hfil$\mapstochar\mkern5mu$\hfil\cr$\to$\cr}}}


\newcommand{\subres}{\trianglelefteq}
\newcommand{\subrctx}[2]{{#1} \subres {#2}}
\makeatletter
\newcommand{\subst}[2][]{\ext@arrow 0359\Rightarrowfill@{#1}{#2}}
\makeatother


\newenvironment{eqns}{\begin{array}{r@{\hspace{0.3em}}c@{\hspace{0.3em}}l}}{\end{array}}


\newcommand{\mat}[1]{\mathbf{#1}}
\newcommand{\vct}[1]{\mathbf{#1}}
\DeclarePairedDelimiter\basis{\langle}{\rvert}


\newcommand{\name}{\ensuremath{\lambda \resctxperm R}}


\DeclareMathOperator\kit{Kit}
\newcommand{\kitrel}{\mathbin{\blacklozenge}}


\DeclareMathOperator\id{id}
\DeclareMathOperator\inl{inl}
\DeclareMathOperator\inr{inr}
\DeclareMathOperator\Idx{Idx}


\newcommand{\zero}{\ensuremath{\rescomment 0}}
\newcommand{\linear}{\ensuremath{\rescomment 1}}
\newcommand{\unrestricted}{\ensuremath{\rescomment \omega}}
\newcommand{\instDILL}{\rescomment{01\omega}}

\newcommand{\unused}{\ensuremath{\rescomment{\mathit{unused}}}}
\newcommand{\true}{\ensuremath{\rescomment{\mathit{true}}}}
\newcommand{\valid}{\ensuremath{\rescomment{\mathit{valid}}}}
\newcommand{\instPD}{\ensuremath{\rescomment{\mathit{utv}}}}

\newcommand{\aval}[1]{\rescomment{\mathrm{#1}}}

% Beamer hacks

\makeatletter
\newcommand{\subst}[2][]{\ext@arrow 0359\Rightarrowfill@{#1}{#2}}
\makeatother

\title{Linear metatheory via linear algebra}
\author{Bob Atkey\inst{1} \and James Wood\inst{1}}
\institute{\inst{1}University of Strathclyde}
\date{SPLS, June 2019}

\begin{document}
\setlength{\abovedisplayskip}{0pt}
\frame{\titlepage}
\begin{frame}
  \frametitle{La logique lin\'eaire}
  \begin{itemize}
  \item A logic of \emph{resources} or \emph{interaction},
    rather than truth or provability.
    \pause
  \item Each variable is used exactly once.
    \pause
  \item Contexts are multisets of formulae.
    \[
      \inferrule*[right=$\otimes$-I]
      {\Gamma \vdash M : A \\ \Delta \vdash N : B}
      {\Gamma, \Delta \vdash (M, N) : A \otimes B}
    \]
    \pause
  \item We should be able to replace variables by terms (substitution).
    \pause
  \item $\oc A$ is like $A$, but can be discarded and duplicated.
    \pause
  \item Substitution is really difficult!
    \[
      \inferrule*[right=$\oc$-I]
      {\Gamma \vdash M : A}
      {\oc\Gamma \vdash [M] : \oc A}
    \]
  \end{itemize}
\end{frame}
\begin{frame}
  \frametitle{A substitute for linear logic}
  \includegraphics[trim={15cm 15cm 18cm 9cm},clip,width=\textwidth]{bread.jpeg}
  \begin{itemize}
  \item Type theory is ever evolving.
  \item More general, more well-behaved.
    \pause
  \item BREAD
  \end{itemize}
\end{frame}
\begin{frame}
  \frametitle{So how you gonna do it?}
  \begin{itemize}
  \item<1-> $\ctx{\Gamma}{\uncover<3->R} \vdash M : A$
  \item<2-> $\ctxvar{x_1}{A_1}{\uncover<3->{\rho_1}}, \ldots,
    \ctxvar{x_n}{A_n}{\uncover<3->{\rho_n}} \vdash M : A$
  \item<4-> The variable rule:
    \[
      \inferrule*[right=Var]
      {\uncover<5->{\rescomment 0, \ldots, \rescomment 0, \rescomment 1,
          \rescomment 0, \ldots, \rescomment 0 \subres \rescomment{\rho_1},
          \ldots, \rescomment{\_}, \rescomment{\pi},
          \rescomment{\_}, \ldots, \rescomment{\rho_n}}}
      {\ctxvar{x_1}{A_1}{\uncover<5->{\rho_1}}, \ldots,
        \ctxvar{y}{B}{\uncover<5->{\pi}}, \ldots,
        \ctxvar{x_n}{A_n}{\uncover<5->{\rho_n}} \vdash y : B}
    \]
  \item<6-> Weakening admissible at $\rescomment 0$ demand:
    \[
      \inferrule*[right=Weak]
      {\ctx{\Gamma}{P} \vdash M : A
        \\ \rescomment{\vec 0} \subres \resctx Q
      }
      {\ctx{\Gamma}{P}, \ctx{\Delta}{Q} \vdash M : A}
    \]
  \item<7-> Examples:
    \setbeamercovered{transparent}
    \begin{enumerate}
    \item<8-> $\ctxvar{x}{A}{1}, \ctxvar{y}{B}{0} \vdash x : A$
      (linearity; DILL from Bar96)
    \item<9-> $\ctxvar{x}{\mathbb Z}{\uparrow}, \ctxvar{y}{\mathbb Z}{\downarrow}
      \vdash x - y : \mathbb Z$ (monotonicity; Arntzenius)
    \item<10-> $\ctxvar{x}{\mathrm{Secret}}{\mathrm{prv}}
      \nvdash \_ : \excl{\mathrm{pub}}{\mathrm{Secret}}$
      (privacy; DCC from ABHR99)
    \end{enumerate}
  \item<7-> Other related work: PD99; RP10; POM14; GS14; BGMZ14;
    the Granule project
  \end{itemize}
\end{frame}
\setbeamercovered{invisible}
\begin{frame}
  \frametitle{Annotation posemiring}
  \begin{itemize}
  \item \rescomment{Annotations} form a partially ordered semiring.
  \item A \emph{partially ordered semiring}
    $(\mathscr R, \subres, 0, +, 1, \cdot)$ is:
    \begin{itemize}
    \item A \emph{partial order} $(\mathscr R, \subres)$
    \item A \emph{semiring} $(\mathscr R, 0, +, 1, \cdot)$, which is:
      \begin{itemize}
      \item A \emph{commutative monoid} $(\mathscr R , 0, +)$
      \item A \emph{monoid} $(\mathscr R , 1, \cdot)$
      \item Such that $\cdot$ distributes over $0$ and $+$ on both sides
      \end{itemize}
    \item Such that $+$ and $\cdot$ are both monotonic in both arguments
    \end{itemize}
    \pause
  \item Ur-example:
    $(\rescomment{\mathbb N} \cup \{\rescomment\omega\}, \sqsupseteq,
    \rescomment 0, +, \rescomment 1, \times)$ ``how many''
    \begin{tikzpicture}[scale=0.75]
      \node (omega) at (0,2) {$\rescomment\omega$};
      \node (zero) at (-4,0) {$\rescomment 0$};
      \node (one) at (-2,0) {$\rescomment 1$};
      \node (mid) at (0,0) {$\cdots$};
      \node (n) at (2,0) {$\rescomment n$};
      \node (end) at (4,0) {$\cdots$};
      \draw (zero) -- (omega)
            (one) -- (omega)
            (mid) -- (omega)
            (n) -- (omega)
            (end) -- (omega);
    \end{tikzpicture}
  \end{itemize}
\end{frame}
%\begin{frame}
%  \frametitle{Non-use}
%  \begin{displaymath}
%    \inferrule*[right=Var]
%    {\Gamma_x = A
%      \\ \begin{pmatrix}
%        \rescomment{\vec 0} & \rescomment 1x & \rescomment{\vec 0}
%      \end{pmatrix} \subres \resctx R
%    }
%    {\ctx{\Gamma}{R} \vdash x : A}
%  \end{displaymath}
%  \begin{displaymath}
%    \inferrule*[right=Weak]
%    {\ctx{\Gamma}{P} \vdash M : A
%      \\ \rescomment{\vec 0} \subres \resctx Q
%    }
%    {\ctx{\Gamma}{P}, \ctx{\Delta}{Q} \vdash M : A}
%  \end{displaymath}
%\end{frame}
\setbeamercovered{transparent}
\begin{frame}
  \frametitle{What are we dealing with?}
  \begin{mathpar}
    \onslide<1,5>{
      % Sharing
      \inferrule*[right=$\&$-I]
      {\ctx{\Gamma}{R} \vdash M : A
        \\\\ \ctx{\Gamma}{R} \vdash N : B
      }
      {\ctx{\Gamma}{R} \vdash \wth{M}{N} : \withT{A}{B}}
    }
    \and
    \onslide<2,5>{
      % Splitting
      \inferrule*[right=$\otimes$-I]
      {\ctx{\Gamma}{P} \vdash M : A
        \\\\ \ctx{\Gamma}{Q} \vdash N : B
        \\\\ \subrctx{\resctx P + \resctx Q}{\resctx R}
      }
      {\ctx{\Gamma}{R} \vdash \ten{M}{N} : \tensor{A}{B}}
    }
    \and
    \onslide<3,5>{
      % Scaling
      \inferrule*[right=$\excl{\rho}{}$-I]
      {\ctx{\Gamma}{P} \vdash M : A
        \\\\ \subrctx{\rescomment \rho \resctx P}{\resctx R}
      }
      {\ctx{\Gamma}{R} \vdash \bang{M} : \excl{\rho}{A}}
    }
    \and
    \onslide<4,5>{
      % Binding
      %\inferrule*[right=$\excl{\rho}{}$-E]
      %{\ctx{\Gamma}{P} \vdash M : \excl{\rho}{A}
      %  \\\\ \ctx{\Gamma}{Q}, \ctxvar{x}{A}{\rho} \vdash N : B
      %  \\\\ \subrctx{\resctx P + \resctx Q}{\resctx R}
      %}
      %{\ctx{\Gamma}{R} \vdash \bm{B}{M}{x}{N} : B}
      \inferrule*[right=$\to$-I]
      {\ctx{\Gamma}{R}, \ctxvar{x}{A}{\rho} \vdash M : B}
      {\ctx{\Gamma}{R} \vdash \lam{x}{M} : \rescomment\rho A \to B}
    }
  \end{mathpar}
\end{frame}
\setbeamercovered{invisible}
\begin{frame}
  \frametitle{Simultaneous substitution}
    \[
      \inferrule*[lab=Subst]
      {\ctx{\Gamma}{\uncover<2->{P}} \vdash M : A
        \\ {\begin{pmatrix}
            \phantom{S_1}\rescomment{\mathllap{\only<2-3>{\resctx Q}}\mathllap{\only<4>{\wn}}\mathllap{\only<5->{S_1}}}\Delta
            \vdash N_1 : \Gamma_1 \\
            \vdots \\
            \phantom{S_m}\rescomment{\mathllap{\only<2-3>{\resctx Q}}\mathllap{\only<4>{\wn}}\mathllap{\only<5->{S_m}}}\Delta
            %\rescomment{\uncover<2->{\Alt<-3>{\resctx Q}{\Alt<4>{\wn}{S_m}}}}\Delta
            \vdash N_m : \Gamma_m
          \end{pmatrix}}}
      {\ctx{\Delta}{\uncover<2->{Q}} \vdash M\{\vec N\} : A}
    \]
  \begin{itemize}
  \item<3->
    \begin{minipage}[t]{\linewidth}
      \begin{overprint}
        \onslide<3>
        Wrong! Try identity substitution on \\
        $\ctx{\Gamma}{P} = \ctx{\Delta}{Q} = \ctxvar{x}{B}{1}, \ctxvar{y}{C}{1}$.

        \onslide<4->
        We need to split $\resctx{Q}$ up.
      \end{overprint}
      \[
        \begin{pmatrix}
          \ctxvar{x}{B}{1}, \ctxvar{y}{C}{\alt<3>{1}{0}} \alt<3>\nvdash\vdash x
          : B \\
          \ctxvar{x}{B}{\alt<3>{1}{0}}, \ctxvar{y}{C}{1} \alt<3>\nvdash\vdash y
          : C
        \end{pmatrix}
      \]
    \end{minipage}
  \item<5->
    What is the relationship between $\resctx P$, $\resctx Q$, and
    $\rescomment{S}$?
  \end{itemize}
\end{frame}
\begin{frame}
  \frametitle{Substitution, stripped down}
    \[
      \inferrule*[lab=Subst]
      {\alt<-2>{\resctx P}{\rescomment{\pi_1}x_1, \ldots, \rescomment{\pi_m}x_m}\only<1>\Gamma \vdash M \only<1>{: A}
        \\ {\begin{pmatrix}
            \alt<-2>{\rescomment{S_1}}{\rescomment{\sigma_{1,1}}y_1, \ldots, \rescomment{\sigma_{1,n}}y_n}\only<1>\Delta \vdash N_1 \only<1>{: \Gamma_1} \\
            \vdots \\
            \alt<-2>{\rescomment{S_m}}{\rescomment{\sigma_{m,1}}y_1, \ldots, \rescomment{\sigma_{m,n}}y_n}\only<1>\Delta \vdash N_m \only<1>{: \Gamma_m} \\
          \end{pmatrix}}}
      {\alt<-2>{\resctx Q}{\rescomment{\rho_1}y_1, \ldots, \rescomment{\rho_n}y_n}\only<1>\Delta \vdash M\{\vec N\} \only<1>{: A}}
    \]
  \begin{itemize}
  \item
    What is the relationship between $\resctx P$, $\resctx Q$, and
    $\rescomment{S}$?
  \item<4->
    $\rescomment{\rho_j}$ is a weighted sum of $\rescomment{\sigma_{-,j}}$
    according to $\rescomment{\pi_{-}}$.
  \item<5->
    \[
      \begin{pmatrix}
        \rescomment{\pi_1} & \cdots & \rescomment{\pi_m}
      \end{pmatrix}
      \begin{pmatrix}
        \rescomment{\sigma_{1,1}} & \cdots & \rescomment{\sigma_{1,n}} \\
        \vdots & \ddots & \vdots \\
        \rescomment{\sigma_{m,1}} & \cdots & \rescomment{\sigma_{m,n}} \\
      \end{pmatrix}
      \subres
      \begin{pmatrix}
        \rescomment{\rho_1} & \cdots & \rescomment{\rho_n}
      \end{pmatrix}
    \]
  \end{itemize}
\end{frame}
\begin{frame}
  \frametitle{What is a substitution?}
    \[
      \inferrule*[lab=Subst]
      {\rescomment{\pi_1}x_1, \ldots, \rescomment{\pi_m}x_m \vdash M
        \\ {\begin{pmatrix}
            \rescomment{\sigma_{1,1}}y_1, \ldots, \rescomment{\sigma_{1,n}}y_n \vdash N_1 \\
            \vdots \\
            \rescomment{\sigma_{m,1}}y_1, \ldots, \rescomment{\sigma_{m,n}}y_n \vdash N_m \\
          \end{pmatrix}}}
      {\rescomment{\rho_1}y_1, \ldots, \rescomment{\rho_n}y_n \vdash M\{\vec N\}}
    \]
  \begin{itemize}
  \item Notation: $\langle x \rvert$ stands for the basis vector at $x$.
  \item A simultaneous substitution
    $\ctx{\Delta}{Q} \subst{} \ctx{\Gamma}{P}$ is:
    \begin{itemize}
    \item A $\left| \Gamma \right|$-by-$\left| \Delta \right|$ matrix
      $\rescomment S$, such that
    \item For each $x : A$ in $\Gamma$, a term $N_x$ \\ such that
      $(\langle x \rvert\rescomment S)\Delta \vdash N_x : A$
    \item $\resctx P \rescomment S \subres \resctx Q$
    \end{itemize}
  \end{itemize}
\end{frame}
\setbeamercovered{transparent}
\begin{frame}
  \frametitle{Substitution on \TirName{$\&$-I}}
  \[
    \inferrule*[right=Subst]{
      \inferrule*[right=$\&$-I]
      {\ctx{\Gamma}{P} \vdash M : A
        \\ \ctx{\Gamma}{P} \vdash N : B
      }
      {\ctx{\Gamma}{P} \vdash \wth{M}{N} : \withT{A}{B}}
      \\
      \ctx{\Delta}{Q} \subst{\sigma} \ctx{\Gamma}{P}
    }
    {\ctx{\Delta}{Q} \vdash \wth{M\{\sigma\}}{N\{\sigma\}} : \withT{A}{B}}
  \]
  \pause
  \[
    \inferrule*[right=$\&$-I]
    {
      \inferrule*[right=Subst]
      {
        \ctx{\Gamma}{P} \vdash M : A
        \\\\
        \ctx{\Delta}{Q} \subst{\sigma} \ctx{\Gamma}{P}
      }
      {\ctx{\Delta}{Q} \vdash M\{\sigma\} : A}
      \\
      \inferrule*[Right=Subst]
      {
        \ctx{\Gamma}{P} \vdash N : B
        \\\\
        \ctx{\Delta}{Q} \subst{\sigma} \ctx{\Gamma}{P}
      }
      {\ctx{\Delta}{Q} \vdash N\{\sigma\} : B}
    }
    {\ctx{\Delta}{Q} \vdash \wth{M\{\sigma\}}{N\{\sigma\}} : \withT{A}{B}}
  \]
\end{frame}
\begin{frame}
  \frametitle{Substitution on \TirName{$\otimes$-I}}
  \[
    \inferrule*[right=Subst]{
      \inferrule*[right=$\otimes$-I]
      {\ctx{\Gamma}{P} \vdash M : A
        \\ \ctx{\Gamma}{Q} \vdash N : B
        \\\\ \subrctx{\resctx P + \resctx Q}{\resctx R}
      }
      {\ctx{\Gamma}{R} \vdash \ten{M}{N} : \tensor{A}{B}}
      \\
      \ctx{\Delta}{R'} \subst{\sigma} \ctx{\Gamma}{R}
    }
    {\ctx{\Delta}{R'} \vdash \ten{M\{\sigma\}}{N\{\sigma\}} : \tensor{A}{B}}
  \]
  \pause
  \begin{displaymath}
    \begin{matrix}
      \resctx{P'} := \resctx P \rescomment S &
      \resctx{Q'} := \resctx Q \rescomment S \\
      \ctx{\Delta}{P'} \subst{\sigma} \ctx{\Gamma}{P} &
      \ctx{\Delta}{Q'} \subst{\sigma} \ctx{\Gamma}{Q}
    \end{matrix}
  \end{displaymath}
  \pause
  \begin{displaymath}
    \resctx{P'} + \resctx{Q'}
    \uncover<1-2,4->{= \resctx P \rescomment S + \resctx Q \rescomment S}
    \uncover<1-2,5->{= (\resctx P + \resctx Q) \rescomment S}
    \uncover<1-2,6->{\subres \resctx R \rescomment S}
    \uncover<1-2,3,7>{\subres \resctx{R'}}
  \end{displaymath}
  \[
    \inferrule*[right=$\otimes$-I]
    {\ctx{\Delta}{P'} \vdash M\{\sigma\} : A
      \\ \ctx{\Delta}{Q'} \vdash N\{\sigma\} : B
      \\\\ \subrctx{\resctx{P'} + \resctx{Q'}}{\resctx{R'}}
    }
    {\ctx{\Delta}{R'} \vdash \ten{M\{\sigma\}}{N\{\sigma\}} : \tensor{A}{B}}
  \]
\end{frame}
\begin{frame}
  \frametitle{Substitution on \TirName{$\oc_{\rescomment\rho}$-I}}
  \[
    \inferrule*[right=Subst]{
      \inferrule*[right=$\oc_{\rescomment\rho}$-I]
      {\ctx{\Gamma}{P} \vdash M : A
        \\ \subrctx{\rescomment \rho \resctx P}{\resctx R}
      }
      {\ctx{\Gamma}{R} \vdash \bang{M} : \excl{\rho}{A}}
      \\
      \ctx{\Delta}{R'} \subst{\sigma} \ctx{\Gamma}{R}
    }
    {\ctx{\Delta}{R'} \vdash \bang{M\{\sigma\}} : \excl{\rho}{A}}
  \]
  \pause
  \begin{displaymath}
    \begin{matrix}
      \resctx{P'} := \resctx P \rescomment S \\
      \ctx{\Delta}{P'} \subst{\sigma} \ctx{\Gamma}{P}
    \end{matrix}
  \end{displaymath}
  \pause
  \begin{displaymath}
    \rescomment \rho \resctx{P'}
    = \rescomment \rho (\resctx P \rescomment S)
    = (\rescomment \rho \resctx P) \rescomment S
    \subres \resctx R \rescomment S
    \subres \resctx{R'}
  \end{displaymath}
  \[
    \inferrule*[right=$\oc_{\rescomment\rho}$-I]
    {\ctx{\Delta}{P'} \vdash M\{\sigma\} : A
      \\ \subrctx{\rescomment \rho \resctx{P'}}{\resctx{R'}}
    }
    {\ctx{\Delta}{R'} \vdash \bang{M\{\sigma\}} : \excl{\rho}{A}}
  \]
\end{frame}
\begin{frame}
  \frametitle{Substitution on \TirName{$\to$-I}}
  \[
    \inferrule*[right=Subst]{
      \inferrule*[right=$\to$-I]
      {\ctx{\Gamma}{R}, \ctxvar{x}{A}{\rho} \vdash M : B}
      {\ctx{\Gamma}{R} \vdash \lam{x}{M} : \rescomment\rho A \to B}
      \\
      \ctx{\Delta}{R'} \subst{\sigma} \ctx{\Gamma}{R}
    }
    {\ctx{\Delta}{R'} \vdash \lam{x}{M\{\sigma,x \mapsto x\}}
      : \rescomment\rho A \to B}
  \]
  \pause
  \begin{mathpar}
    \rescomment{S'} := \begin{pmatrix}
      & & & \rescomment 0 \\
      & \rescomment S & & \vdots \\
      & & & \rescomment 0 \\
      \rescomment 0 & \cdots & \rescomment 0 & \rescomment 1
    \end{pmatrix}
    \and
    \begin{pmatrix}\resctx R & \rescomment\rho\end{pmatrix}\rescomment{S'}
    = \begin{pmatrix}\resctx{R'} & \rescomment\rho\end{pmatrix}
  \end{mathpar}
  \pause
  \[
    \forall(z : C) \in \Gamma, x : A.~\exists N_z.~
    (\langle z \rvert\rescomment{S'})\Delta \vdash N_z : C
  \]
  \pause
  Cases:
  \begin{itemize}
  \item
    $
    \forall(z : C) \in \Gamma.~\exists N_z.~
    (\langle z \rvert\rescomment{S}) \Delta, \ctxvar{x}{A}{0} \vdash N_z : C
    $
  \item
    $
    \begin{pmatrix}\rescomment 0 & \cdots & \rescomment 0\end{pmatrix} \Delta,
    \ctxvar{x}{A}{1} \vdash x : A
    $
  \end{itemize}
\end{frame}
\begin{frame}
  \frametitle{DILL}
  \begin{mathpar}
    \begin{array}{c|ccc}
      + & \aval 0 & \aval 1 & \aval\omega \\ \hline
      \aval 0 & \aval 0 & \aval 1 & \aval\omega \\
      \aval 1 & \aval 1 & \aval\omega & \aval\omega \\
      \aval\omega & \aval\omega & \aval\omega & \aval\omega
    \end{array}
    \and
    \begin{array}{c|ccc}
      \cdot & \aval 0 & \aval 1 & \aval\omega \\ \hline
      \aval 0 & \aval 0 & \aval 0 & \aval 0 \\
      \aval 1 & \aval 0 & \aval 1 & \aval\omega \\
      \aval\omega & \aval 0 & \aval\omega & \aval\omega
    \end{array}
    \and
    \begin{tikzpicture}[scale=0.75,baseline={(0,1)}]
      \node (omega) at (0,2) {$\aval\omega$};
      \node (zero) at (-2,0) {$\aval 0$};
      \node (one) at (2,0) {$\aval 1$};
      \draw (zero) -- (omega)
            (one) -- (omega);
    \end{tikzpicture}
    \and
    \Gamma; \Delta \vdash M : A \leftrightsquigarrow
    \vec{\rescomment\omega}\Gamma, \vec{\rescomment 1}\Delta,
    \vec{\rescomment 0}\Theta \vdash M : A
  \end{mathpar}
\end{frame}
\begin{frame}
  \frametitle{Pfenning, Davies}
  \begin{mathpar}
    0 = \aval{unused}
    \and
    1 = \aval{true}
    \\
    \begin{array}{c|ccc}
      + & \aval u & \aval t & \aval v \\ \hline
      \aval u & \aval u & \aval t & \aval v \\
      \aval t & \aval t & \aval t & \aval v \\
      \aval v & \aval v & \aval v & \aval v
    \end{array}
    \and
    \begin{array}{c|ccc}
      \cdot & \aval u & \aval t & \aval v \\ \hline
      \aval u & \aval u & \aval u & \aval u \\
      \aval t & \aval u & \aval t & \aval v \\
      \aval v & \aval u & \aval v & \aval v
    \end{array}
    \and
    \begin{tikzpicture}[scale=0.75,baseline={(0,0)}]
      \node (valid) at (0,2) {$\aval{valid}$};
      \node (true) at (0,0) {$\aval{true}$};
      \node (unused) at (0,-2) {$\aval{unused}$};
      \draw (true) -- (valid)
            (unused) -- (true);
    \end{tikzpicture}
  \end{mathpar}
\end{frame}
\begin{frame}
  \frametitle{Conclusions}
  \begin{itemize}
  \item A notion of substitution that remains constant, even though contexts
    change.
  \item Lesson: don't remove things from scope; mark them depleted
    ($\rescomment 0$).
  \item Agda code at \url{https://github.com/laMudri/quantitative}
  \item Translation table:
    \begin{tabular}{|l|p{170pt}|}
      \hline
      Syntax & Linear algebra \\
      \hline
      Annotation context & Vector \\
      \TirName{Var} annotation & Basis vector \\
      Weakening & Embedding into a space of higher dimension \\
      A substitution & A linear map \\
      \hline
    \end{tabular}
  \end{itemize}
\end{frame}

\appendix

\begin{frame}
  \frametitle{Substitution on \TirName{Var}}
  \[
    \inferrule*[right=Subst]{
      \inferrule*[right=Var]
      {(x : A) \in \Gamma \\ \langle x \rvert \subres \resctx R}
      {\ctx{\Gamma}{R} \vdash x : A}
      \\
      \ctx{\Delta}{R'} \subst{\sigma} \ctx{\Gamma}{R}
    }
    {\ctx{\Delta}{R'} \vdash \sigma_x : A}
  \]
  \begin{displaymath}
    \langle x \rvert\rescomment S
    \subres \resctx R \rescomment S
    \subres \resctx{R'}
  \end{displaymath}
  From $\sigma$, we have
  $(\langle x \rvert\rescomment S)\Delta \vdash \sigma_x : A$.
\end{frame}
\begin{frame}
  \frametitle{Single substitution}
  \[
    \inferrule*[right=SingleSubst]{
      \ctx{\Gamma}{P}, \ctxvar{x}{A}{\rho} \vdash M : B
      \\ \ctx{\Gamma}{Q} \vdash N : A
    }
    {(\resctx P + \rescomment\rho\resctx Q)\Gamma \vdash M[N/x]}
  \]
  \[
    \resctx P + \rescomment\rho\resctx Q \subst{[N/x]} \resctx P, \rescomment\rho x
  \]
  Need $\begin{pmatrix}\resctx P & \rescomment\rho\end{pmatrix}\rescomment S
  \subres \resctx P + \rescomment\rho\resctx Q$
  \[
    S := \begin{pmatrix}
      \rescomment 1 & \rescomment 0 & \cdots & \rescomment 0 \\
      \rescomment 0 & \rescomment 1 & \cdots & \rescomment 0 \\
      \vdots & \vdots & \ddots & \vdots \\
      \rescomment 0 & \rescomment 0 & \cdots & \rescomment 1 \\
      \resctx Q_1 & \resctx Q_2 & \cdots & \resctx Q_n
    \end{pmatrix}
  \]
\end{frame}

\end{document}
